\documentclass[11pt,oneside,openany]{book}

\usepackage[T1]{fontenc}
\usepackage{lmodern}
\usepackage[letterpaper,margin=1.10in,includeheadfoot]{geometry}
\usepackage{amsmath,amssymb,amsthm,mathtools}
\usepackage{microtype}
\usepackage{xcolor}
\usepackage{titlesec}
\usepackage{fancyhdr}
\usepackage{needspace}
\usepackage{hyperref}

\ifdefined\pdftrailerid
  \pdftrailerid{MathGenome Advanced Field Theory}
\fi

\definecolor{ink}{HTML}{172923}
\definecolor{deepgreen}{HTML}{2B5242}
\definecolor{rulegreen}{HTML}{8A9F92}
\hypersetup{
  colorlinks=true,linkcolor=deepgreen,urlcolor=deepgreen,citecolor=deepgreen,
  pdftitle={Advanced Field Theory},
  pdfsubject={Infinite Galois Theory},
  pdfauthor={MathGenome},bookmarksnumbered=true
}
\setcounter{tocdepth}{1}
\setcounter{secnumdepth}{2}
\setlength{\headheight}{15pt}
\setlength{\emergencystretch}{2em}

\titleformat{\chapter}[display]
  {\normalfont\color{ink}}{\color{deepgreen}\large\scshape Chapter \thechapter}{9pt}
  {\huge\bfseries}
\titlespacing*{\chapter}{0pt}{4pt}{25pt}
\titleformat{\section}{\normalfont\Large\bfseries\color{ink}}{\thesection}{0.7em}{}
\titleformat{\subsection}{\normalfont\large\bfseries\color{ink}}{\thesubsection}{0.7em}{}

\pagestyle{fancy}
\fancyhf{}
\fancyhead[L]{\small\sffamily\color{deepgreen}MathGenome}
\fancyhead[R]{\small\nouppercase{\leftmark}}
\fancyfoot[C]{\small\thepage}
\renewcommand{\headrulewidth}{0.35pt}
\renewcommand{\chaptermark}[1]{\markboth{\thechapter.\ #1}{}}
\fancypagestyle{plain}{%
  \fancyhf{}\fancyfoot[C]{\small\thepage}\renewcommand{\headrulewidth}{0pt}}

\newtheorem{theorem}{Theorem}[chapter]
\newtheorem{lemma}[theorem]{Lemma}
\newtheorem{proposition}[theorem]{Proposition}
\newtheorem{corollary}[theorem]{Corollary}
\theoremstyle{definition}
\newtheorem{definition}[theorem]{Definition}
\newtheorem{example}[theorem]{Example}
\theoremstyle{remark}
\newtheorem{remark}[theorem]{Remark}
\allowdisplaybreaks[1]

\begin{document}
\frontmatter
\hypersetup{pageanchor=false}
\begin{titlepage}
\thispagestyle{empty}
\thispdfpagelabel{Cover}
\begin{flushleft}
\vspace*{0.65in}
{\large\scshape\color{deepgreen}MathGenome\par}
\vspace{0.35in}
{\color{rulegreen}\rule{0.94\textwidth}{0.65pt}\par}
\vspace{0.43in}
{\color{ink}\fontsize{36}{42}\selectfont\bfseries
Advanced\\[4pt] Field Theory\par}
\vspace{0.32in}
{\color{deepgreen}\Large Infinite Galois Theory\par}
\vfill
{\large Mathematical lecture notes\par}
\vspace{0.10in}
{\normalsize MathGenome v0.21.0 dependency baseline\par}
\vspace*{0.45in}
\end{flushleft}
\end{titlepage}
\hypersetup{pageanchor=true}

\tableofcontents

\chapter*{Scope and conventions}
\addcontentsline{toc}{chapter}{Scope and conventions}

The finite Galois correspondence describes an extension through the
subgroups of its automorphism group. For an arbitrary algebraic Galois
extension, the information about any particular element still sits in a
finite Galois subextension. The Krull topology records simultaneous
agreement on these finite parts. We use it to identify the full group
with an inverse limit and to determine precisely which subgroups are
recovered from their fixed fields. A faithful compact group action in
the opposite direction constructs a Galois extension, and in particular
realizes any profinite group over a suitable field.

Throughout, a Galois extension is algebraic, normal, and separable.
In the general argument \(L/K\) is Galois and
\(G=\operatorname{Gal}(L/K)\). For an intermediate field \(E\) write
\(G_E=\operatorname{Gal}(L/E)\) for its pointwise stabilizer.
If \(E/K\) is Galois, write \(Q_E=\operatorname{Gal}(E/K)\) and
\(r_E:G\to Q_E\) for restriction. For finite Galois stages
\(E\subseteq E'\) the transition map \(r_{E',E}\) goes from
\(Q_{E'}\) to \(Q_E\). The finite groups carry the discrete topology.
We use left cosets in \(G/G_E\) when discussing arbitrary \(E\).

The first chapter recalls the finite field-extension tools from
\cite[Chapters 12--14]{d07} and proves the extension and restriction
facts needed in the infinite case. The next chapter constructs the
topology and glues compatible automorphisms. The third treats closed
subgroups, open subgroups and closed normal quotients. The last computes
the absolute Galois group of a finite field, tracking the arithmetic
Frobenius \(x\mapsto x^q\) through every restriction. The source map
records the references and proof dependencies of the numbered results,
together with the relevant interfaces at the fixed Lean baseline.

\mainmatter
\chapter[Finite subextensions]{From Finite Subextensions to an Algebraic Galois Extension}
\label{ch:finite-subextensions}

Fix a field extension \(L/K\). We identify \(K\) with its image in
\(L\). An intermediate field is a subfield of this particular \(L\)
containing \(K\), and its scalar structures are the inclusions. If an
abstract extension of \(K\) is given with a \(K\)-embedding into \(L\),
we use its image as an intermediate field and transport its operations
through that embedding. For \(a\) algebraic over a field \(F\), write
\(\mu_{a,F}\) for the monic minimal polynomial of \(a\) over \(F\).

\section{Finite results and passage to an intermediate field}

\begin{proposition}[Finite ingredients]\label{res:F01}
Let \(a\) be algebraic over \(K\), and let \(T\) be an extension of
\(K\). The \(K\)-embeddings \(K(a)\to T\) correspond to the roots of
\(\mu_{a,K}\) in \(T\), by evaluation at \(a\). The field generated by
all the roots of a separable polynomial in a splitting extension is
finite Galois over \(K\).

If \(E/K\) is finite Galois, then
\[
 M\longmapsto\operatorname{Gal}(E/M),\qquad
 H\longmapsto E^H
\]
give inverse inclusion-reversing maps between the intermediate fields
of \(E/K\) and the subgroups of \(\operatorname{Gal}(E/K)\).
In particular \(E^{\operatorname{Gal}(E/K)}=K\).
For an intermediate field \(M\) Galois over \(K\), restriction
\(\operatorname{Gal}(E/K)\to\operatorname{Gal}(M/K)\) is surjective
with kernel \(\operatorname{Gal}(E/M)\).
\end{proposition}

These are the finite results used below. For the first statement the
presentation \(K(a)=K[X]/(\mu_{a,K})\) explains the root parametrization
and shows that the embedding associated with a given root is unique.
The splitting field facts, the finite correspondence, and the finite
restriction sequence are established in \cite[Chapters 12 and 14]{d07}.
See also \cite[Chapters 2 and 3]{milne} and
\cite[Tag 09DW]{stacks}. When counting embeddings of a finite
separable extension \(M/K\) into an extension containing the conjugates
of all elements of \(M\), we use a primitive element \(a\) of \(M/K\).
Its minimal polynomial then has \([M:K]\) distinct roots in that
extension, and the first statement gives the same number of embeddings.

We call \(L/K\) Galois if it is algebraic, normal, and separable.
Normality means that every irreducible polynomial in \(K[X]\) having
a root in \(L\) splits in \(L[X]\). In an algebraic extension separability
can be checked on the minimal polynomial of every element. We write
\[
 G=\operatorname{Gal}(L/K),\qquad G_M=\operatorname{Gal}(L/M)
\]
for the group of \(K\)-automorphisms and the subgroup fixing an
intermediate field \(M\) pointwise. If \(M/K\) is Galois, put
\(Q_M=\operatorname{Gal}(M/K)\). For a subgroup \(H\) of \(G\), \(L^H\)
denotes the field \(\{x\in L:h(x)=x\ \text{for every }h\in H\}\).

\begin{proposition}[The upper part of a Galois tower]\label{res:F02}
If \(L/K\) is Galois and \(M\) is an intermediate field, then \(M/K\)
is algebraic and separable, and \(L/M\) is Galois.
\end{proposition}
\begin{proof}
For \(a\in M\), the polynomial \(\mu_{a,K}\) is the same whether \(a\)
is viewed in \(M\) or in \(L\). This proves algebraicity and separability
for \(M/K\). For every \(b\in L\), the polynomial \(\mu_{b,K}\),
regarded as a polynomial over \(M\), annihilates \(b\). Thus \(b\)
is algebraic over \(M\), and \(\mu_{b,M}\) divides \(\mu_{b,K}\) in
\(M[X]\). Over \(L\) the latter is a product of distinct linear factors.
Each of its monic divisors over \(L\) therefore has the same splitting
and distinctness properties. This proves normality and separability
of \(L/M\). Compare \cite[Proposition 7.3]{milne} and
\cite[Tag 09I2]{stacks}.
\end{proof}

\section{Finite Galois subextensions}

\begin{proposition}[A finite Galois subextension containing a finite set]
\label{res:F03}
If \(L/K\) is Galois, every finite subset \(S\subseteq L\) lies in an
intermediate field \(E\) which is finite Galois over \(K\).
\end{proposition}
\begin{proof}
The empty subset lies in \(K\). Otherwise, take the product of the
distinct polynomials \(\mu_{a,K}\) for \(a\in S\). Each factor is
separable and splits in \(L\), and two distinct monic irreducible
factors have disjoint root sets. The product is therefore separable.
Inside \(L\) take the field \(E\) generated by all of its roots.
This field contains \(S\), and it is precisely the splitting field
of the product inside \(L\). Proposition~\ref{res:F01} makes it finite
Galois over \(K\). The choice of distinct factors also applies if
different elements of \(S\) are conjugate. See \cite[Chapter 7, p.~95]{milne}
and \cite[Tag 0BMG]{stacks}.
\end{proof}

\begin{theorem}[The directed union of the finite Galois subextensions]
\label{res:F04}
In a Galois extension \(L/K\), the finite Galois intermediate fields,
ordered by inclusion, form a nonempty directed set. The compositum of
two of them inside \(L\) is again one of them, and their union is \(L\).
Conversely, if an extension of \(K\) is the union of a nonempty directed
family of finite Galois intermediate fields over \(K\), the extension
is Galois.
\end{theorem}
\begin{proof}
The ground field \(K\) is a finite Galois intermediate field.
Let \(E_1,E_2\) be two such fields. In each \(E_i\) choose a finite
set of field generators over \(K\), for example a \(K\)-basis.
The minimal polynomials of these generators over \(K\) split in
\(E_i\). The field generated by their roots in \(L\) is contained
in \(E_i\) and contains all its generators, so it equals \(E_i\).
Thus \(E_1E_2\) is the field generated by the roots of the polynomials
from both lists. It is the splitting field of their product with
duplicate monic irreducible factors omitted. That product is
separable, and Proposition~\ref{res:F01} makes \(E_1E_2/K\) finite
Galois. This provides a common upper bound for \(E_1\) and \(E_2\).
Proposition~\ref{res:F03} applied to singletons proves the union
assertion.

For the converse, directedness ensures that two elements of the
union lie in a common field, so the union is a field. Let \(a\)
lie in the union. One field in the family contains \(a\), and
its being finite over \(K\) proves that
\(a\) is algebraic over \(K\). The minimal polynomial \(\mu_{a,K}\)
is separable and splits in that same field. This holds for every
\(a\) and proves both separability and normality of the union over
\(K\). See \cite[Proposition 3.21, Chapter 7, and Example 7.26]{milne},
\cite[Theorem 2.11.1]{rz}, and \cite[Tag 0BU2]{stacks}.
\end{proof}

\section{Extension of embeddings and restriction}

\begin{theorem}[Extension of an intermediate-field embedding]
\label{res:F05}
Suppose \(L/K\) is Galois. If \(M\) is an intermediate field and
\(\iota:M\to L\) is a \(K\)-embedding, there is a \(K\)-automorphism
\(\sigma:L\to L\) with \(\sigma|_M=\iota\).
\end{theorem}
\begin{proof}
Choose an algebraic closure \(\Omega\) of \(L\), using the given
inclusion \(L\subseteq\Omega\). Consider the pairs \((N,\tau)\) with
\(M\subseteq N\subseteq L\) and \(\tau:N\to\Omega\) a \(K\)-embedding
extending \(\iota\) as a map into \(\Omega\). Order them by extension
of both domain and map. They include \((M,\iota)\). Along a chain
the union of the domains is a field and the embeddings agree on
overlaps. Their union is an embedding, since any pair of arguments
lies in a common member of the chain. Hence every nonempty chain
has an upper bound, and Zorn's lemma gives a maximal pair.

If its domain \(N\) is proper in \(L\), choose \(x\in L\setminus N\).
By algebraicity over \(K\), \(x\) is algebraic over \(N\).
Transport \(\mu_{x,N}\) along \(\tau\); the result is irreducible
over \(\tau(N)\) and has a root in \(\Omega\).
The quotient presentation of a simple extension in
Proposition~\ref{res:F01} extends \(\tau\) to \(N(x)\) by sending
\(x\) to that root. This contradicts maximality. We obtain a
\(K\)-embedding \(\tau:L\to\Omega\).

For \(x\in L\), its image \(\tau(x)\) is a root of \(\mu_{x,K}\).
Normality of \(L/K\) puts all the roots of this polynomial in \(L\),
so we can regard \(\tau\) as a \(K\)-endomorphism of \(L\).
It is surjective: the finite set of roots of \(\mu_{x,K}\) is
stable under \(\tau\), and the restriction to this set is injective,
hence a permutation. In particular \(x\) has a preimage in \(L\).
This is true for every \(x\), so \(\tau\) is the required
automorphism. See \cite[Theorem 6.6 and Proposition 7.4]{milne},
\cite[Chapter 12]{d07}, and \cite[Tag 0BME]{stacks}.
\end{proof}

\begin{corollary}[Stable intermediate fields]\label{res:F06}
In a Galois extension \(L/K\), an intermediate field \(M\) is Galois
over \(K\) if and only if \(\sigma(M)=M\) for every \(\sigma\in G\).
\end{corollary}
\begin{proof}
If \(M/K\) is Galois and \(a\in M\), then \(\sigma(a)\) is a root of
\(\mu_{a,K}\), whose roots all lie in \(M\). Thus \(\sigma(M)\subseteq M\).
Apply this also to \(\sigma^{-1}\) to get equality.
Conversely, suppose \(M\) is stable under \(G\). For any \(a\in M\)
every root \(b\) of \(\mu_{a,K}\) in an algebraic closure of \(L\)
already lies in \(L\). By Proposition~\ref{res:F01}, \(a\mapsto b\)
defines a \(K\)-embedding \(K(a)\to L\).
Theorem~\ref{res:F05} extends it to a member of \(G\).
Stability of \(M\) gives \(b\in M\), so \(\mu_{a,K}\) splits in \(M\).
This proves normality over \(K\). Algebraicity and separability
follow from Proposition~\ref{res:F02}. See
\cite[Corollary 7.5 and Theorem 7.13(d)]{milne}.
\end{proof}

\begin{proposition}[Surjective restriction and composition]
\label{res:F07}
Suppose \(L/K\) is Galois. If \(M/K\) is a Galois intermediate
extension, restriction gives a surjective homomorphism
\[
 r_M:G\longrightarrow Q_M,\qquad r_M(\sigma)=\sigma|_M,\qquad
 \ker r_M=G_M.
\]
For any two Galois intermediate fields \(E\subseteq E'\),
restriction also gives a surjective homomorphism
\(r_{E',E}:Q_{E'}\to Q_E\).
These maps satisfy \(r_{E,E}=\operatorname{id}_{Q_E}\),
\(r_E=r_{E',E}\circ r_{E'}\), and
\[
 r_{E'',E}=r_{E',E}\circ r_{E'',E'}
 \quad\text{for } E\subseteq E'\subseteq E''
\]
among Galois intermediate fields.
\end{proposition}
\begin{proof}
Corollary~\ref{res:F06} ensures that \(\sigma|_M\) maps \(M\) onto
itself, so restriction takes values in \(Q_M\).
The restriction of a composite agrees with the composite of the
restrictions, since every restricted map has domain and image \(M\).
An element of \(Q_M\), composed with the inclusion of \(M\) into \(L\),
is a \(K\)-embedding \(M\to L\), and Theorem~\ref{res:F05} extends
it to a member of \(G\). This proves surjectivity. The kernel is
precisely the set of maps acting identically on \(M\).
For \(r_{E',E}\) apply the same argument with \(E'\) as top field.
Each identity of restriction maps follows by evaluating both
sides on an element of the smallest field involved. In particular
these identities make the finite-stage restriction maps an inverse
system. The continuity assertions for restriction will use the
topology constructed in the next chapter. See
\cite[Proposition 7.7 and Example 7.26]{milne} and
\cite[Tag 0BMK]{stacks}.
\end{proof}

\chapter{Krull Topology and the Inverse Limit}\label{ch:krull-topology}

Let \(L/K\) be a Galois extension, and put
\(G=\operatorname{Gal}(L/K)\). Products in \(G\) are compositions,
with the right factor applied first. For an intermediate field \(E\),
write \(G_E=\operatorname{Gal}(L/E)\), viewed as a subgroup of \(G\).
For a finite subset \(S\subseteq L\), put
\[
 G(S)=\{\sigma\in G:\sigma(x)=x\text{ for all }x\in S\}.
\]
In particular, \(G(S)=G_{K(S)}\). Fixing \(S\) fixes every expression
obtained from its elements and \(K\) by the field operations, and hence
fixes the field they generate.

\section{The neighborhood basis and restriction}

\begin{proposition}[Finite pointwise stabilizers]\label{res:T01}
There is a unique topology with basis consisting of the left cosets
\(\sigma G(S)\), where \(S\) ranges over the finite subsets of \(L\).
In this topology \(G\) is a topological group, and the groups \(G(S)\)
are an open neighborhood basis at the identity. This is the
\emph{Krull topology}.
\end{proposition}

\begin{proof}
The subsets \(G(S)\) are subgroups, with
\[
 G(S)\cap G(S')=G(S\cup S'),\qquad
 \sigma G(S)\sigma^{-1}=G(\sigma S).
\]
For the second identity, \(\sigma h\sigma^{-1}\) fixes each \(\sigma x\)
if and only if \(h\) fixes each \(x\). Every point belongs to a proposed
basic open set. If \(\gamma\) belongs to two basic cosets, each coset can
be written using \(\gamma\) as representative. Their intersection then
equals \(\gamma G(S\cup S')\). This proves the basis property.
Equivalently, \(\sigma G(S)\) consists of all automorphisms which agree
with \(\sigma\) on \(S\).

For continuity of multiplication, fix a basic neighborhood
\(\sigma\tau G(S)\) at a product. If \(\tau'\) agrees with \(\tau\) on
\(S\) and \(\sigma'\) agrees with \(\sigma\) on \(\tau S\), then
\(\sigma'\tau'(x)=\sigma\tau(x)\) for \(x\in S\). Thus the product of
the neighborhoods \(\sigma G(\tau S)\) and \(\tau G(S)\) is contained
in the given neighborhood. For inversion, if \(\sigma'\) agrees
with \(\sigma\) on \(\sigma^{-1}S\), then
\(\sigma'(\sigma^{-1}x)=x\) for \(x\in S\). Applying the inverse of
\(\sigma'\) proves that \((\sigma')^{-1}\) agrees with \(\sigma^{-1}\)
on \(S\). This supplies the required basic neighborhood at \(\sigma\).
In a topological group translations are homeomorphisms, so the identity
basis determines the entire topology. See
\cite[Section 7, following 7.2, and Proposition 7.6]{milne}.
\end{proof}

Let \(\mathcal E\) be the set of finite Galois intermediate fields of
\(L/K\). For \(E\in\mathcal E\), put
\[
 Q_E=\operatorname{Gal}(E/K),\qquad
 r_E\colon G\longrightarrow Q_E,\quad \sigma\longmapsto \sigma|_E,
 \qquad N_E=\ker r_E=G_E.
\]
We regard every \(Q_E\) as a discrete group. Theorem~\ref{res:F04}
gives the two indexing properties used here: \(\mathcal E\) is
directed by inclusion, and its union is \(L\).

\begin{proposition}[A basis of finite quotient kernels]\label{res:T02}
The subgroups \(N_E\), for \(E\in\mathcal E\), form an open neighborhood
basis at the identity. Each is clopen and normal of finite index, and
\(\bigcap_{E\in\mathcal E}N_E=\{1\}\). If \(L/K\) is finite, the Krull
topology is discrete.
\end{proposition}

\begin{proof}
For each \(E\in\mathcal E\), choose a \(K\)-basis of \(E\) and let \(S_E\)
be its finite set of elements. Then \(N_E=G(S_E)\), since a
\(K\)-automorphism fixing the basis pointwise fixes \(E\). Conversely,
Proposition~\ref{res:F03} places any finite \(S\subseteq L\) in some
\(E\in\mathcal E\). For that \(E\) one has \(N_E\subseteq G(S)\), giving
the claimed basis. Proposition~\ref{res:F07} identifies \(N_E\) as the
kernel of the surjection \(r_E\), so it is normal and its index is
\(|Q_E|\). Every open subgroup has open cosets. The complement of
\(N_E\) is the union of its other cosets, making \(N_E\) closed as well.
An automorphism in all \(N_E\) fixes their fields' union \(L\) pointwise.
If \(L/K\) is finite we may take \(E=L\), for which \(N_E=\{1\}\).
Translations then give open singletons throughout \(G\).
Compare \cite[Propositions 7.6--7.8]{milne} and
\cite[Section 2.11]{rz}.
\end{proof}

\begin{proposition}[Initial topologies and continuous action]\label{res:T03}
The Krull topology is the initial topology of the evaluations
\(\sigma\mapsto\sigma(x)\) to discrete \(L\), for all \(x\in L\). It is
also the initial topology of the restrictions \(r_E\) to discrete
\(Q_E\), for all \(E\in\mathcal E\). The action \(G\times L\to L\) is
jointly continuous when \(L\) is discrete. If \(M/K\) is a Galois
intermediate extension, then restriction
\(G\to\operatorname{Gal}(M/K)\) is continuous for the Krull topologies.
\end{proposition}

\begin{proof}
The nonempty inverse image of a singleton under the evaluation at \(x\)
is a coset of \(G(\{x\})\), and any \(\sigma G(S)\) is the intersection
of the corresponding fibers for \(x\in S\). These identities give the
first topology. Each nonempty fiber of \(r_E\) is a coset of \(N_E\).
Conversely all these cosets are fibers. By Proposition~\ref{res:T02}
they form a basis, giving the second initial topology.

At \((\sigma,x)\), the action is constant on
\(\sigma G(\{x\})\times\{x\}\). This is an open neighborhood, which
proves joint continuity. It also yields the corresponding universal
property for an action map: if a map \(u\colon X\to G\) from a
topological space has continuous induced map
\((z,x)\mapsto u(z)(x)\) on \(X\times L\), then each
\(z\mapsto u(z)(x)\) is continuous by fixing \(x\), so \(u\) is
continuous by the first initial description.

For the restriction to \(M\), Proposition~\ref{res:F07} gives the
homomorphism, since \(M/K\) is Galois. For \(x\in M\), its composite
with evaluation at \(x\) is the evaluation in \(L\) with values in the
discrete subset \(M\). This composite is continuous as a map into
discrete \(M\), because the inverse images of its singletons are
the same as the corresponding inverse images for evaluation in \(L\).
Applying the first initial description to \(M/K\) gives continuity
of restriction, including when \(M/K\) is infinite.
See \cite[Propositions 7.6--7.7]{milne} and
\cite[Tags 0BMJ and 0BMK]{stacks}.
\end{proof}

\section{An inverse system of finite groups}

\begin{proposition}[The compatible-tuple limit]\label{res:T04}
For \(E\subseteq E'\) in \(\mathcal E\), let
\(r_{E',E}\colon Q_{E'}\to Q_E\) be restriction. These onto maps
form an inverse system. With coordinatewise group operations and
the subspace topology from \(\prod_{E\in\mathcal E}Q_E\), the group
\[
 P=\left\{(\sigma_E)_{E\in\mathcal E} :
       r_{E',E}(\sigma_{E'})=\sigma_E
       \text{ whenever }E\subseteq E'\right\}
\]
is its inverse limit in topological groups.
\end{proposition}

\begin{proof}
The restrictions make sense because \(E/K\) is Galois, and
Proposition~\ref{res:F07}, applied to \(E'/K\), gives surjectivity.
Restrictions are the identity when the two fields agree, and
\(r_{E'',E}=r_{E',E}r_{E'',E'}\) for \(E\subseteq E'\subseteq E''\).
The maps are continuous between discrete groups. Every compatibility
equation is preserved by multiplication and inversion because its
transition map is a homomorphism. Thus \(P\) is a subgroup of the
product, where multiplication and inversion are continuous by the
product topology.

Let \(X\) be a topological group with continuous homomorphisms
\(f_E\colon X\to Q_E\) satisfying \(r_{E',E}f_{E'}=f_E\).
Their componentwise pairing takes values in \(P\). It is a continuous
homomorphism to \(P\) by the product and subspace topologies, with all
the prescribed components. An element of \(P\) is its tuple of
components, so these conditions determine the factorization uniquely.
This is the stated limit property.
See \cite[Examples 7.22--7.23 and 7.26]{milne}.
\end{proof}

\begin{theorem}[Gluing compatible automorphisms]\label{res:T05}
The homomorphism
\[
 \rho\colon G\longrightarrow P,\qquad
 \sigma\longmapsto (\sigma|_E)_{E\in\mathcal E}
\]
is an isomorphism of groups. Its composite with projection to \(Q_E\)
is \(r_E\), and each such projection is surjective.
\end{theorem}

\begin{proof}
The finite restrictions of an automorphism are compatible by their
composition law. Two automorphisms with the same restrictions agree
on the union \(L\), proving injectivity.

Given \((\sigma_E)\in P\) and \(x\in L\), choose \(E\in\mathcal E\)
containing \(x\) and set \(\sigma(x)=\sigma_E(x)\). If \(E'\) also
contains \(x\), choose a common upper field \(E''\) by directedness.
The two proposed values equal \(\sigma_{E''}(x)\) by compatibility.
To test addition or multiplication, choose one field containing
both inputs. It also contains their sum and product, and the
corresponding \(\sigma_E\) respects those operations. The glued map
fixes \(K\) and preserves \(1\) for the same reason.

The tuple \((\sigma_E^{-1})\) is compatible because restrictions are
homomorphisms. Its glued map is a two-sided inverse to \(\sigma\):
on any \(E\), both maps preserve \(E\) and their restrictions are
inverse automorphisms. Thus \(\sigma\in G\) and \(\rho(\sigma)\) is
the prescribed tuple. The identification of projections follows
from the definition of \(\rho\), and their surjectivity follows
from Proposition~\ref{res:F07}.
See \cite[Example 7.26]{milne},
\cite[Theorem 2.11.1]{rz}, and \cite[Tag 0BU2]{stacks}.
\end{proof}

\begin{theorem}[Topological inverse limit and profiniteness]\label{res:T06}
The isomorphism \(\rho\) is a homeomorphism. Consequently \(G\) is
compact, Hausdorff, and totally disconnected, and is profinite.
\end{theorem}

\begin{proof}
Continuity of every \(r_E\) makes \(\rho\) continuous into the product
and hence into \(P\). Its inverse is continuous because every
\(r_E\) composed with that inverse is the coordinate projection
on \(P\), and the \(r_E\) generate the Krull topology by
Proposition~\ref{res:T03}.

Each \(Q_E\) is finite discrete, hence compact and Hausdorff.
Their product is compact by the Tychonoff product theorem proved in
\cite[Chapter~5, \texttt{d10:05:tychonoff}]{d10}. It is Hausdorff
because distinct tuples differ in a coordinate.
For each \(E\subseteq E'\) the compatibility equation defining \(P\)
is the inverse image of the diagonal of \(Q_E\times Q_E\) under
\[
  (\sigma_D)_D \longmapsto
      \bigl(r_{E',E}(\sigma_{E'}),\sigma_E\bigr).
\]
This map is continuous and the diagonal is closed. Taking the
intersection of these closed sets shows that \(P\) is a closed
subspace of the compact Hausdorff product. If two points of \(P\)
differ, a fiber of one coordinate separates them by a clopen
subset of \(P\). No connected subset can contain both points.
Thus \(P\) and \(G\) are totally disconnected. The presentation
by finite discrete groups gives profiniteness directly.
See \cite[Propositions 7.8 and 7.25 and Example 7.26]{milne}
and \cite[Theorem 2.11.1]{rz}.
\end{proof}

\section{Closure, cofinality, and a change of base}

\begin{proposition}[Closure on finite levels]\label{res:T07}
For any subset \(A\subseteq G\),
\[
 \overline A=\bigcap_{E\in\mathcal E}r_E^{-1}(r_E(A)).
\]
If \(\mathcal E_0\subseteq\mathcal E\) is cofinal for inclusion,
then it suffices to use \(E\in\mathcal E_0\) in this formula.
For the same cofinal subset the coordinate restriction
\(P\to\varprojlim_{D\in\mathcal E_0}Q_D\) is an isomorphism of
topological groups.
\end{proposition}

\begin{proof}
For each \(g\in G\), its basic neighborhood \(gN_E\) is
\(r_E^{-1}(\{r_E(g)\})\). It meets \(A\) if and only if
\(r_E(g)\in r_E(A)\). Membership in \(\overline A\) is equivalent
to this meeting condition for every basic neighborhood, proving the
formula. If \(E\subseteq D\), then \(N_D\subseteq N_E\), so the
cofinal subgroups still furnish a neighborhood basis. The same
proof gives the restricted closure formula.

To see the limit comparison, start with a compatible tuple over
\(\mathcal E_0\). For any \(E\in\mathcal E\), choose
\(D\in\mathcal E_0\) with \(E\subseteq D\) and define its \(E\)-entry
to be \(r_{D,E}(\sigma_D)\). If two choices \(D,D'\) are made, first
take a common upper bound in \(\mathcal E\) and then an upper
bound for it in \(\mathcal E_0\). The compatibility equations
through this upper field make the two entries equal. For
\(E\subseteq E'\), choose \(D\) containing \(E'\), so the transition
law also makes the reconstructed entries compatible. This proves
that reconstruction is inverse to coordinate restriction.
It is continuous since each reconstructed entry is the composite
of a projection on \(\mathcal E_0\) and a transition. Both maps
respect the group operations coordinatewise. Compare
\cite[Proposition 7.12(b) and Example 7.26]{milne}.
\end{proof}

\begin{proposition}[The Krull topology over an intermediate field]
\label{res:T08}
For any intermediate field \(M\) of \(L/K\), the topology on \(G_M\)
induced from \(G\) equals its intrinsic Krull topology for \(L/M\).
\end{proposition}

\begin{proof}
The extension \(L/M\) is Galois by Proposition~\ref{res:F02}. Its
Krull topology has identity neighborhoods given by the pointwise
stabilizers of finite \(S\subseteq L\) in \(G_M\), by
Proposition~\ref{res:T01} over the new base. These are precisely
\(G_M\cap G(S)\). The intersections also form the induced
neighborhood basis at the identity, and their translates by
points of \(G_M\) therefore give both topologies. This comparison
uses the same elements of the top field \(L\) on the two sides and
applies to every intermediate field \(M\).
See \cite[Theorem 2.11.3]{rz}.
\end{proof}

\chapter{The Infinite Galois Correspondence}
\label{ch:galois-correspondence}

Throughout Sections~\ref{sec:fixed-fields}--\ref{sec:normal-quotients},
let \(L/K\) be Galois, let \(G=\operatorname{Gal}(L/K)\) carry the Krull
topology, and let \(G_M=\operatorname{Gal}(L/M)\) for an intermediate
field \(M\). For a subgroup \(H\leq G\) we write
\[
  L^H=\{x\in L : h(x)=x\text{ for every }h\in H\}.
\]
This is a field between \(K\) and \(L\), since every \(h\) preserves
addition, multiplication, and inverses of nonzero elements. We keep
the notation \(Q_E=\operatorname{Gal}(E/K)\) and
\(r_E:G\to Q_E\) for a finite Galois intermediate field \(E\).
The finite fixed-field theorem applied in \(E\) will control both
the restriction of \(L^H\) to \(E\) and the closure of \(H\) in \(G\).

\section{Fixed fields and closure}
\label{sec:fixed-fields}

\begin{proposition}[The fixed field of the full group]
\label{res:C01}
For a Galois extension \(L/K\), the full group
\(G=\operatorname{Gal}(L/K)\) has fixed field \(K\). More precisely,
if \(x\in L\setminus K\), then there is a finite Galois intermediate
field \(E\) containing \(x\) and an element of \(Q_E\) moving \(x\)
that is the restriction of an element of \(G\).
\end{proposition}

\begin{proof}
Choose \(E\) by \ref{res:F03}. The finite fixed-field theorem
of \ref{res:F01} gives \(E^{Q_E}=K\), so some \(\tau\in Q_E\)
satisfies \(\tau(x)\ne x\). Surjectivity of \(r_E\) in \ref{res:F07}
provides \(g\in G\) with \(g|_E=\tau\). Therefore \(g(x)\ne x\).
All elements of \(K\) are fixed by the definition of \(G\).
\cite[Proposition 7.9]{milne}
\end{proof}

\begin{proposition}[Recovery of an intermediate field]
\label{res:C02}
For every intermediate field \(M\), the subgroup \(G_M\) is closed
and \(L^{G_M}=M\).
\end{proposition}

\begin{proof}
For \(x\in M\), the stabilizer of \(x\) in \(G\) is the inverse
image of \(\{x\}\) under evaluation at \(x\). Evaluation is continuous
to the discrete space \(L\) by \ref{res:T03}, so the stabilizer is
closed. Their intersection over \(x\in M\) is \(G_M\). By \ref{res:F02}
the extension \(L/M\) is Galois, and \ref{res:C01} with base field \(M\)
gives \(L^{G_M}=M\).
\cite[Proposition 7.12(a)]{milne}
\end{proof}

\begin{lemma}[Fixed fields on a finite layer]
\label{res:C03}
For any subgroup \(H\leq G\) and any finite Galois intermediate field
\(E\), one has
\[
  L^H\cap E=E^{r_E(H)}
  \qquad\text{and}\qquad
  r_E(H)=\operatorname{Gal}\bigl(E/(L^H\cap E)\bigr).
\]
\end{lemma}

\begin{proof}
Each element \(h\in H\) preserves \(E\) because \(E/K\) is normal.
For \(x\in E\) the condition \(h(x)=x\) for every \(h\in H\)
is therefore the condition that each restricted automorphism in
\(r_E(H)\) fix \(x\). This proves the first equality. The image
\(r_E(H)\) is a subgroup of the finite Galois group \(Q_E\).
Apply the finite correspondence of \ref{res:F01} to that subgroup
and the intermediate field \(E^{r_E(H)}\) to obtain the second equality.
\cite[Tag 0BML]{stacks}
\end{proof}

\begin{theorem}[Closure and the infinite Galois correspondence]
\label{res:C04}
For any subgroup \(H\leq G\),
\[
    G_{L^H}=\overline H,
    \qquad L^{\overline H}=L^H .
\]
Consequently the assignments \(M\mapsto G_M\) and \(H\mapsto L^H\)
give inverse inclusion-reversing bijections between the intermediate
fields of \(L/K\) and the closed subgroups of \(G\).
\end{theorem}

\begin{proof}
The group \(G_{L^H}\) is closed by \ref{res:C02} and contains \(H\),
so it contains \(\overline H\). If \(g\in G_{L^H}\), then \(r_E(g)\)
fixes every element of \(L^H\cap E\). By \ref{res:C03} this says that
\(r_E(g)\in r_E(H)\). This is true for every finite Galois intermediate
field \(E\), so the finite-projection closure criterion \ref{res:T07}
gives \(g\in\overline H\). We have proved the first equality. In
particular \(\overline H\) is a subgroup, and \ref{res:C02} applied
to \(M=L^H\) gives the second equality.

For a closed subgroup \(H\), the first equality now reads
\(G_{L^H}=H\). This and \ref{res:C02} prove that the assignments
are inverse. Enlarging an intermediate field imposes more pointwise
fixing conditions on its subgroup. Enlarging a subgroup imposes
more conditions on its fixed field. Thus the two inverses reverse
inclusion.
\cite[Proposition 7.12 and Theorem 7.13(a)]{milne}
\cite[Theorem 2.11.3]{rz}
\cite[Tag 0BML]{stacks}
\end{proof}

\begin{corollary}[Lattice operations and conjugation]
\label{res:C05}
For a family of intermediate fields \((M_i)_{i\in I}\), let
\(\bigvee_i M_i\) be their compositum in \(L\), and put \(H_i=G_{M_i}\).
Then
\[
  G_{\bigvee_i M_i}=\bigcap_i H_i,
  \qquad
  G_{\bigcap_i M_i}
     =\overline{\left\langle\bigcup_i H_i\right\rangle}.
\]
An empty compositum is taken to be \(K\), and an empty intersection
of intermediate fields is \(L\). For \(g\in G\) one also has
\[
  G_{g(M)}=gG_Mg^{-1},
  \qquad
  L^{gHg^{-1}}=g(L^H)\quad\text{for every }H\leq G.
\]
\end{corollary}

\begin{proof}
Fixing all \(M_i\) pointwise is equivalent to fixing the field they
generate pointwise, which proves the first lattice equality. Write
\(J=\langle\bigcup_i H_i\rangle\). An element of \(L\) is fixed by \(J\)
exactly when it is fixed by all the \(H_i\), because a product of
automorphisms fixing an element and their inverses still fixes the
element. Thus \(L^J=\bigcap_i L^{H_i}=\bigcap_i M_i\) by
\ref{res:C02}. Apply \ref{res:C04} to \(J\).

An automorphism \(a\in G\) fixes all elements \(g(m)\), \(m\in M\),
exactly when \(g^{-1}ag\) fixes every \(m\in M\). This proves the
first conjugation formula. For the second, write a candidate element
as \(g(x)\). Since \(ghg^{-1}(g(x))=g(h(x))\), it is fixed by all
the conjugates of \(H\) exactly when \(x\) is fixed by \(H\).
\cite[Theorem 7.13(c) and Remark 7.14]{milne}
\end{proof}

\section{Finite intermediate fields and cosets}

\begin{theorem}[Open subgroups and finite extensions]
\label{res:C06}
For an intermediate field \(M\), the subgroup \(G_M\) is open if
and only if \(M/K\) is finite. More generally, a closed subgroup
of \(G\) is open if and only if it has finite index.
\end{theorem}

\begin{proof}
If \(M/K\) is finite, a finite \(K\)-basis \(S\) for \(M\) generates
\(M\) as a field over \(K\). The subgroup \(G_M\) is then the
pointwise stabilizer of \(S\), which is open by \ref{res:T01}.
Conversely, if \(G_M\) is open, it contains the subgroup \(N_E=G_E\)
for some finite Galois intermediate field \(E\), by the neighborhood
basis of \ref{res:T02}. Taking fixed fields and applying \ref{res:C02}
gives \(M=L^{G_M}\subseteq L^{G_E}=E\). The inclusion makes \(M\)
a \(K\)-subspace of the finite-dimensional space \(E\), so \(M/K\)
is finite.

If \(H\leq G\) is open, its left cosets are open and form a disjoint
cover of the compact space \(G\) by \ref{res:T06}. A finite subcover
must contain all cosets, so the index is finite. If \(H\) is closed
of finite index, its complement is the finite union of its other
left cosets. Each is closed because left translation is a
homeomorphism. Thus \(H\) is open.
\cite[Theorem 7.13(b)]{milne}
\cite[Sections 2.1 and 2.11]{rz}
\end{proof}

\begin{proposition}[Cosets as embeddings]
\label{res:C07}
The restriction of automorphisms induces a bijection of sets
\[
  G/G_M \longrightarrow \operatorname{Emb}_K(M,L),
  \qquad gG_M\longmapsto g|_M ,
\]
where \(G/G_M\) denotes the set of left cosets. In particular, if
\(M/K\) is finite, then
\[
    [G:G_M]=[M:K].
\]
\end{proposition}

\begin{proof}
The equality \(g|_M=g'|_M\) holds if and only if
\(g^{-1}g'\in G_M\), so the map is well-defined and injective.
Every \(K\)-embedding \(M\to L\) extends to an element of \(G\)
by \ref{res:F05}, giving surjectivity. For finite \(M/K\), it is
separable by \ref{res:F02}. Choose a primitive element \(x\) by
the finite preliminaries \ref{res:F01}. Its minimal polynomial
over \(K\) has degree \([M:K]\), has distinct roots, and splits
in \(L\) by normality. Each root gives exactly one \(K\)-embedding
of \(K(x)=M\) into \(L\), proving the count. This identifies the
degree with the coset index for every finite intermediate field,
including those for which \(G_M\) is not normal.
\cite[Proposition 7.4 and Theorem 7.13(b)]{milne}
\cite[Chapters 12 and 14]{d07}
\end{proof}

\section{Normal subgroups and topological quotients}
\label{sec:normal-quotients}

\begin{corollary}[Normality in the correspondence]
\label{res:C08}
Let \(H\) be a closed subgroup of \(G\) with fixed field \(M\).
Then \(H\) is normal in \(G\) if and only if \(M/K\) is Galois.
The open normal subgroups correspond precisely to the finite
Galois intermediate extensions of \(K\).
\end{corollary}

\begin{proof}
By \ref{res:C04}, \(H=G_M\). For every \(g\in G\), the conjugate
\(gG_Mg^{-1}\) equals \(G_{g(M)}\) by \ref{res:C05}.
By injectivity of the field-to-subgroup correspondence,
this equals \(G_M\) exactly when \(g(M)=M\). By \ref{res:F06},
stability under all \(g\) is equivalent to the Galois condition
on \(M/K\). An open subgroup is closed, since its complement is a
union of open cosets, so the openness condition can be added
using \ref{res:C06}.
\cite[Theorem 7.13(c)--(d)]{milne}
\cite[Tag 0BML]{stacks}
\end{proof}

\begin{theorem}[The profinite exact sequence]
\label{res:C09}
Suppose \(M/K\) is Galois. The inclusion and restriction maps form
a short exact sequence of profinite topological groups
\[
  1\longrightarrow G_M\longrightarrow G
   \xrightarrow{\ r_M\ }\operatorname{Gal}(M/K)
   \longrightarrow 1 .
\]
In particular the induced map
\[
    G/G_M \longrightarrow \operatorname{Gal}(M/K),
    \qquad gG_M\longmapsto g|_M ,
\]
is a topological group isomorphism for the quotient topology on
the left and the Krull topology on the right.
\end{theorem}

\begin{proof}
The subgroup \(G_M\) is closed and normal by \ref{res:C02}
and \ref{res:C08}. Its own Krull topology as \(\operatorname{Gal}(L/M)\)
is exactly the topology induced by its inclusion into \(G\), by
\ref{res:T08}. The restriction \(r_M\) is continuous by \ref{res:T03},
and \ref{res:F07} identifies its kernel as \(G_M\) and proves
surjectivity. Hence it factors as a bijective homomorphism through
the group quotient. The factor is continuous by the definition
of the quotient topology. Its domain is compact as a continuous
image of \(G\), and \(\operatorname{Gal}(M/K)\) is Hausdorff by
\ref{res:T06}. A continuous bijection from a compact space to
a Hausdorff space is a homeomorphism, since it sends closed sets
to compact and hence closed sets. The inclusion is a topological
embedding, and the restriction is a quotient map by this homeomorphism.
All three groups are profinite by \ref{res:T06}, applied to
the Galois extensions \(L/M\), \(L/K\), and \(M/K\).
\cite[Theorem 7.13(d)]{milne}
\cite[Theorem 2.11.3]{rz}
\cite[Tag 0BMM]{stacks}
\end{proof}

\section{Restriction after a base change}

\begin{theorem}[Galois base change in a common overfield]
\label{res:C10}
Let \(E\) and \(M\) be extensions of \(K\) in a common overfield,
and suppose \(E/K\) is Galois. Then \(EM/M\) and \(E/(E\cap M)\)
are Galois, and restriction is an isomorphism of topological groups
\[
  \operatorname{Gal}(EM/M)
     \longrightarrow \operatorname{Gal}\bigl(E/(E\cap M)\bigr),
     \qquad g\longmapsto g|_E .
\]
\end{theorem}

\begin{proof}
The extension \(E/(E\cap M)\) is Galois by \ref{res:F02}.
By \ref{res:F04}, \(E\) is the directed union of its finite
Galois intermediate fields \(E_i\) over \(K\). Each \(E_i\)
is the splitting field of a finite collection of separable
polynomials over \(K\): take minimal polynomials of a finite
set of field generators and use normality to put all their
roots in \(E_i\). Over \(M\) the same polynomials remain separable,
and their roots generate \(E_iM\) over \(M\). Hence \(E_iM/M\)
is finite Galois. The union of the \(E_iM\) is a field containing
both \(E\) and \(M\), and it is contained in \(EM\), so it equals
\(EM\). The converse direction of \ref{res:F04} now proves
that \(EM/M\) is Galois.

Every \(g\in\operatorname{Gal}(EM/M)\) is a \(K\)-automorphism and
sends \(E\) into \(E\) by normality of \(E/K\): for \(x\in E\),
the element \(g(x)\) is a root of the \(K\)-minimal polynomial of
\(x\), and all those roots already belong to \(E\). The same applies
to \(g^{-1}\), so \(g|_E\) is an automorphism fixing \(E\cap M\).
The restriction homomorphism is injective because \(E\) and \(M\)
generate \(EM\). For a finite subset \(S\subseteq E\), the inverse
image of its pointwise stabilizer in the target is its pointwise
stabilizer in the source, which is open by \ref{res:T01}.
Thus the homomorphism is continuous at the identity and hence
everywhere.

Let \(H\) denote its image. The source is compact and the target
is Hausdorff by \ref{res:T06}, so \(H\) is closed. An element of
\(E\) is fixed by \(H\) exactly when it is fixed by every
element of \(\operatorname{Gal}(EM/M)\). By \ref{res:C01},
the latter group fixes precisely \(M\) in \(EM\). It follows that
\(E^H=E\cap M\). Apply \ref{res:C04} to the Galois extension
\(E/(E\cap M)\) and its closed subgroup \(H\) to obtain
\(H=\operatorname{Gal}(E/(E\cap M))\). Restriction is now a
continuous bijection from compact to Hausdorff as in \ref{res:C09},
and therefore is a homeomorphism.
\cite[Proposition 7.15]{milne}
\cite[Chapter 14, finite base change]{d07}
\end{proof}

\section{Fixed fields of compact groups and realization}

\begin{theorem}[Fixed fields of faithful compact actions]
\label{res:C11}
Let \(G\) be a compact topological group acting faithfully by field
automorphisms on a field \(E\). Suppose the stabilizer of every
element of \(E\) is open in \(G\). Put \(K=E^G\). Then \(E/K\)
is Galois, and the homomorphism induced by the action
\[
 G \longrightarrow \operatorname{Gal}(E/K)
\]
is an isomorphism of topological groups, with the Krull topology
on the right.
\end{theorem}

\begin{proof}
The orbit of \(a\in E\) is in bijection with the left cosets of its
stabilizer. These cosets are open and cover the compact space \(G\).
A finite subcover therefore accounts for every coset, so every
orbit is finite.

For a finite set \(S\subseteq E\), let \(E_S\) be the field generated
over \(K\) by all the \(G\)-orbits of elements of \(S\). The
generator set is finite and \(G\)-stable. Hence \(G\) acts on
\(E_S\) through a finite group \(H_S\) of field automorphisms:
the action on the generator set determines the restriction to
the whole field. The fixed field of \(H_S\) in \(E_S\) is exactly
\(K\), since an element of \(E_S\) fixed by \(H_S\) is fixed by
all of \(G\). Artin's finite fixed-field theorem, in
\cite[Chapter 14, Artin's fixed-field theorem]{d07} and
\cite[Corollary 3.11]{milne}, now gives that \(E_S/K\) is
finite Galois and that \(H_S=\operatorname{Gal}(E_S/K)\).
In particular the restricted action of \(G\) fills this entire
finite Galois group.

The fields \(E_S\) are directed since \(E_{S\cup T}\) contains
both \(E_S\) and \(E_T\), and their union is \(E\). The converse
in \ref{res:F04} makes \(E/K\) Galois. Any finite Galois
intermediate field is generated by finitely many elements over
\(K\), so it is contained in some \(E_S\).

The homomorphism \(G\to\operatorname{Gal}(E/K)\) is injective
by faithfulness. It is continuous because the inverse image of
a basic pointwise stabilizer on a finite set is an intersection
of finitely many open stabilizers in \(G\), using
\ref{res:T01}. Its image is dense: given an automorphism of
\(E/K\) and any finite set \(S\) on which to match it, take its
restriction to \(E_S\). The equality
\(H_S=\operatorname{Gal}(E_S/K)\) provides an element of \(G\)
with that restriction. This meets each Krull neighborhood of the
given automorphism. The image is also compact, hence closed in
the Hausdorff group of \ref{res:T06}. It is therefore the whole
group. A continuous bijection from compact to Hausdorff is a
homeomorphism, as in \ref{res:C09}.
\cite[Proposition 7.10]{milne}
\end{proof}

\Needspace{16\baselineskip}
\begin{corollary}[Tate's realization theorem]
\label{res:C12}
Every profinite group \(G\) is isomorphic, as a topological
group, to \(\operatorname{Gal}(E/K)\) for some Galois extension
\(E/K\). The fields can be constructed in any specified
characteristic.
\end{corollary}

\begin{proof}
Let \(X\) be the disjoint union of the sets \(G/U\) over the open
normal subgroups \(U\) of \(G\). Such subgroups have finite
index by compactness. The left translation action on \(X\) is
faithful: the kernel on \(G/U\) is \(U\), and the intersection
of the open normal subgroups is trivial. Indeed, in an inverse
limit presentation of \(G\), the finite-coordinate kernels
are open and normal and separate points. Every point of \(X\)
has open stabilizer, because \(U\) fixes each point of the
corresponding component.

Choose a field \(k\) of the desired characteristic and put
\(E=k(t_x:x\in X)\), the rational function field in independent
indeterminates indexed by \(X\). The rule
\(t_x\mapsto t_{g x}\) permutes the variables and extends
to a \(k\)-automorphism for each \(g\in G\). These
automorphisms form a faithful action. A rational function
involves only finitely many variables. Its stabilizer contains
the intersection of their point stabilizers, an open subgroup,
and is therefore open as a union of cosets of that subgroup.
Apply \ref{res:C11} to this action, with \(K=E^G\).
The construction specifies \(K\) as a fixed field of \(E\).
\cite[Proposition 7.27]{milne}
\end{proof}

\chapter{Finite Fields and Arithmetic Frobenius}
\label{ch:finite-fields}

Let \(k\) be a finite field of cardinality \(q=p^f\), with \(p\)
prime, and fix an algebraic closure \(\Omega\) of \(k\).
The finite subfields used here are literal subfields of this one
closure. This choice makes all inclusion and restriction maps part
of the same system.

\section{Finite layers of the algebraic closure}

\begin{proposition}[Root fields in the chosen closure]\label{res:E01}
For each \(n\geq 1\), the roots of \(X^{q^n}-X\) in \(\Omega\)
form a subfield \(k_n\) with \(q^n\) elements. This field contains
\(k\), has \([k_n:k]=n\), and is a finite Galois extension of \(k\).
It is the unique subfield of \(\Omega\) with that cardinality.
\end{proposition}
\begin{proof}
The polynomial splits in \(\Omega\) and its derivative is \(-1\).
All its \(q^n\) roots are therefore distinct. The map
\(a\mapsto a^{q^n}\) is an iterate of the characteristic-\(p\)
Frobenius endomorphism. Its fixed elements contain \(0\) and \(1\)
and are closed under addition, subtraction, multiplication and
inversion of nonzero elements. They form a field. Every \(a\in k\)
satisfies \(a^q=a\), by the order of \(k^\times\) for \(a\ne0\),
so \(k\subseteq k_n\). The cardinality of a finite-dimensional
\(k\)-space of dimension \(d\) is \(q^d\), giving \([k_n:k]=n\).
Since \(k_n\) consists of all roots, it is the splitting field of
this separable polynomial. Finally, in a field of cardinality
\(q^n\) every element satisfies \(a^{q^n}=a\), again by the order of
the multiplicative group. Any such subfield of \(\Omega\) is therefore
contained in \(k_n\), and equality follows by counting.
Compare \cite[Propositions 4.20 and 4.23]{milne},
\cite[Theorem 2.5]{ln}, and \cite[Chapter 13]{d07}.
\end{proof}

\begin{proposition}[The divisor lattice]\label{res:E02}
For positive integers \(m,n\),
\[
\begin{aligned}
   k_m\subseteq k_n &\quad\Longleftrightarrow\quad m\mid n,\\
   k_m\cap k_n &= k_{\gcd(m,n)},&
   k_m k_n &= k_{\operatorname{lcm}(m,n)} .
\end{aligned}
\]
Here \(k_m k_n\) denotes the compositum in \(\Omega\).
The inclusions for \(m\mid n\) are \(k\)-embeddings and satisfy
the identity and composition laws.
\end{proposition}
\begin{proof}
If \(m\mid n\), an element fixed by the \(q^m\)-power map is fixed by
its \(n/m\)-th iterate. In the reverse direction, use the tower formula
for degrees and Proposition~\ref{res:E01}.
The degree of the intersection divides both \(m\) and \(n\). The
intersection contains \(k_{\gcd(m,n)}\), so its degree is their gcd.
The compositum lies in \(k_{\operatorname{lcm}(m,n)}\). Its degree
is a multiple of both \(m\) and \(n\), so equality holds. The embeddings
are inclusions of subfields of one field. Applying two successive
inclusions to an element leaves that element unchanged, which gives
the composition law. See \cite[Corollary 4.21 and Proposition 4.23]{milne}
and \cite[Theorem 2.6]{ln}.
\end{proof}

\begin{corollary}[Exhaustion by finite fields]\label{res:E03}
The \(k_n\) are all the finite intermediate fields of \(\Omega/k\),
and
\[
                \Omega=\bigcup_{n\geq1}k_n=\bigcup_{j\geq1} k_{j!}.
\]
The factorial stages are cofinal in the divisibility order.
\end{corollary}
\begin{proof}
A finite intermediate field of degree \(d\) has \(q^d\) elements, so
by Proposition~\ref{res:E01} it is \(k_d\). For \(a\in\Omega\) the
simple field \(k(a)\) is a finite intermediate field, since \(a\)
is algebraic over \(k\). This proves the first union.
Since \(n\mid n!\), Proposition~\ref{res:E02} proves the second union
and cofinality. The factorial stages are Galois by
Proposition~\ref{res:E01}. Their directed union is Galois by
Theorem~\ref{res:F04}, and its finite Galois subextensions are
exactly the fields in this divisibility system.
See \cite[Propositions 4.23--4.24]{milne}.
\end{proof}

\section{Frobenius as a compatible generator}

\begin{theorem}[Arithmetic Frobenius and its restrictions]\label{res:E04}
The map \(\Phi:\Omega\to\Omega\), \(x\mapsto x^q\), is a
\(k\)-automorphism of infinite order. Write
\(\Phi_n=\Phi|_{k_n}\). Then \(\Phi_n\) has order \(n\) and
generates \(\operatorname{Gal}(k_n/k)\). For \(m\mid n\) the
group \(\operatorname{Gal}(k_n/k_m)\) is generated by
\(\Phi_n^m\) and has order \(n/m\).
\end{theorem}
\begin{proof}
The characteristic-\(p\) Frobenius and hence its \(f\)-th iterate
are field endomorphisms. The latter is injective and fixes \(k\).
It preserves every root field \(k_n\), and on that finite set is
bijective. By Corollary~\ref{res:E03} it is surjective on \(\Omega\).

By construction \(\Phi_n^n=1\). If \(0<d<n\) and
\(\Phi_n^d=1\), every element of \(k_n\) would be a root of
\(X^{q^d}-X\). Its degree is smaller than \(|k_n|\), which is
impossible. Thus the exact order is \(n\), also for \(n=1\)
where the group is trivial. The finite Galois degree formula
gives \(|\operatorname{Gal}(k_n/k)|=n\), proving generation.
A positive power \(\Phi^d\) cannot be the identity on \(\Omega\):
its restriction to \(k_{d+1}\) has nonzero exponent modulo
\(d+1\). If \(m\mid n\), \(\Phi_n^m\) fixes \(k_m\), and in
the cyclic group of order \(n\) its order is \(n/m\). The
extension \(k_n/k_m\) is Galois with degree \(n/m\), by the
tower property and the degree formula. This proves the relative
assertion. See \cite[Proposition 4.20 and Example 7.17]{milne}
and \cite[Theorem 2.21]{ln}.
\end{proof}

\begin{proposition}[Restriction in exponent coordinates]\label{res:E05}
For every \(n\geq1\) the map
\[
     \lambda_n:(\mathbb Z/n\mathbb Z,+)
             \longrightarrow \operatorname{Gal}(k_n/k),
     \qquad \lambda_n(\bar a)=\Phi_n^a
\]
is a group isomorphism. For \(m\mid n\) let
\(\rho_{n,m}:\mathbb Z/n\mathbb Z\to\mathbb Z/m\mathbb Z\)
be reduction. Then
\[
          r_{k_n,k_m}\circ\lambda_n=\lambda_m\circ\rho_{n,m}.
\]
\end{proposition}
\begin{proof}
The group isomorphism follows from the generation and exact order
in Theorem~\ref{res:E04}. The reduction \(\rho_{n,m}\) is well-defined
since any multiple of \(n\) is a multiple of \(m\).
The automorphism \(\Phi_n\) preserves \(k_m\) and restricts to
\(\Phi_m\). Restriction is a homomorphism by Proposition~\ref{res:F07},
so it sends
\(\Phi_n^a\) to \(\Phi_m^a\) for every integer \(a\), including
negative exponents. This proves the displayed equality on every
residue class. In particular the isomorphisms use the same
distinguished generator at every stage. Compare
\cite[Examples 7.17 and 7.26]{milne}.
\end{proof}

\begin{theorem}[The absolute Galois group of a finite field]
\label{res:E06}
Endow
\[
    \widehat{\mathbb Z}
      =\varprojlim_{n}\mathbb Z/n\mathbb Z
\]
with its inverse-limit topology. The index is the set of positive
integers ordered by divisibility. The topology is that of compatible
residue families in the product of the finite discrete groups.
There is a topological group isomorphism
\[
         \Theta:\operatorname{Gal}(\Omega/k)
                          \xrightarrow{\ \sim\ }(\widehat{\mathbb Z},+)
\]
defined by \(\Theta(\sigma)_n=\lambda_n^{-1}(\sigma|_{k_n})\).
It sends \(\Phi\) to the family \(1=(1\bmod n)_n\).
\end{theorem}
\begin{proof}
By Corollary~\ref{res:E03}, the \(k_n\) are exactly the finite
Galois stages of the extension, with order given by
Proposition~\ref{res:E02}. The topological inverse-limit comparison
in Theorem~\ref{res:T06} identifies \(\operatorname{Gal}(\Omega/k)\)
with their compatible systems of automorphisms. The maps
\(\lambda_n^{-1}\) carry compatible systems to compatible residues
by Proposition~\ref{res:E05}. Conversely the \(\lambda_n\) give the
inverse map on compatible residues. Both maps between these limit
spaces are continuous, since every coordinate is given by a continuous
map between finite discrete spaces. The formula for \(\Phi\) follows
from \(\lambda_n(1)=\Phi_n\) for all \(n\).
This is the residue model of the profinite group in
\cite[Example 7.17]{milne} and \cite[Example 2.11.2]{rz}.
\end{proof}

\section{Intermediate fields in the residue model}

\begin{proposition}[A dense proper subgroup]\label{res:E07}
The ordinary cyclic subgroup \(\langle\Phi\rangle\) is dense and
proper in \(\operatorname{Gal}(\Omega/k)\). Its fixed field is \(k\).
\end{proposition}
\begin{proof}
Under \(\Theta\) its image is the group of integer residue families
\((a\bmod n)_n\) with \(a\in\mathbb Z\). Prescribing finitely many
coordinates of a compatible family is achieved by taking an integer
representative at a common multiple of those moduli. Hence this image
is dense.

To see that the image is proper, for each \(n=2^s d\) with \(d\) odd
choose the unique residue class \(b_n\bmod n\) with
\[
                      b_n\equiv0\pmod{2^s},
                \qquad b_n\equiv1\pmod d .
\]
This exists by the Chinese remainder theorem, including moduli
equal to \(1\). For \(m\mid n\), reduction of \(b_n\) satisfies the
corresponding two conditions for \(m\), so it equals \(b_m\).
The family \((b_n)_n\) is compatible. An integer representing it
would be divisible by all powers of \(2\), hence zero, but its
residue modulo \(3\) would be \(1\). Thus it does not come from
an integer. Finally the fixed elements of \(\Phi\) are the roots
of \(X^q-X\), namely \(k\) by Proposition~\ref{res:E01}.
Equivalently a subgroup and its closure have the same fixed field,
as in Theorem~\ref{res:C04}. See \cite[Example 7.17]{milne}
and \cite[Exercise 2.11.4]{rz}.
\end{proof}

\begin{theorem}[Open subgroups and finite intermediate fields]
\label{res:E08}
For each \(n\geq1\), write \(\pi_n:\widehat{\mathbb Z}\to
\mathbb Z/n\mathbb Z\) for projection. Then
\[
                   \ker\pi_n=n\widehat{\mathbb Z}.
\]
Under \(\Theta\) this subgroup corresponds to
\(\operatorname{Gal}(\Omega/k_n)\), has index \(n\), and has
fixed field \(k_n\). All open subgroups arise uniquely this way.
\end{theorem}
\begin{proof}
The image of the integers under \(\pi_n\) is the whole finite
coordinate group, so \(\pi_n\) is surjective and its kernel is
closed and open. It contains \(n\mathbb Z\), viewed in
\(\widehat{\mathbb Z}\) through the residue map. If \(a\in\ker\pi_n\)
and a neighbourhood of \(a\) specifies finitely many moduli,
choose \(N\) divisible by them and by \(n\). An integer
representative of \(a_N\) is a multiple of \(n\) since \(a_n=0\),
and realizes all the prescribed coordinates. Thus
\(\overline{n\mathbb Z}=\ker\pi_n\).

The group \(\widehat{\mathbb Z}\) is compact by the inverse-limit
description in Theorem~\ref{res:E06} and the compactness of
the product of finite discrete groups. Multiplication by \(n\)
is continuous, so its image is compact and thus closed in this
Hausdorff group. Its image contains \(n\mathbb Z\). Conversely, by
the density of \(\mathbb Z\) proved in Proposition~\ref{res:E07},
the image of every point under multiplication by \(n\) belongs
to \(\overline{n\mathbb Z}\). This proves the kernel formula.

Under \(\Theta\), the map \(\pi_n\) reads the restriction to
\(k_n\), expressed in exponent coordinates. Its kernel is
therefore the subgroup fixing \(k_n\), and the index and fixed
field follow from surjectivity and Theorem~\ref{res:C04}.
Every open subgroup has a finite intermediate fixed field by
Theorem~\ref{res:C06}. Corollary~\ref{res:E03} shows this
field is some \(k_n\), uniquely determined by its degree over \(k\).
The closed-subgroup correspondence recovers the subgroup.
Compare \cite[Theorem 7.13(b) and Example 7.17]{milne}
and \cite[Theorem 2.7.2]{rz}.
\end{proof}

\begin{example}[A closed subgroup with infinite index]\label{res:E09}
For a prime \(\ell\) the intermediate field
\[
              M_\ell=\bigcup_{r\geq0}k_{\ell^r}
\]
is infinite over \(k\) and is a proper subfield of \(\Omega\).
Its fixing subgroup corresponds to the closed nonopen subgroup
\[
                       \bigcap_{r\geq0}\ell^r\widehat{\mathbb Z}.
\]
\end{example}
\begin{proof}
The \(k_{\ell^r}\) form an increasing chain by Proposition~\ref{res:E02},
so their union is a field. Its finite subextensions have unbounded
degrees \(\ell^r\), making \([M_\ell:k]\) infinite. Take any prime
\(s\ne\ell\). If \(k_s\subseteq M_\ell\), a finite \(k\)-basis of
\(k_s\) would be contained in one \(k_{\ell^r}\) because the chain
is increasing. This would give \(s\mid\ell^r\), a contradiction.
The subgroup fixing the union is the intersection of the subgroups
fixing each layer. Theorem~\ref{res:E08} gives the displayed
intersection, which is closed. It is nonopen by the infinite degree
and Theorem~\ref{res:C06}, so the same theorem gives infinite index.
Compare
\cite[Remark 7.14 and Example 7.17]{milne}.
\end{proof}


\backmatter
\begin{thebibliography}{9}
\addcontentsline{toc}{chapter}{Bibliography}

\bibitem{milne}
James S.~Milne.
\emph{Fields and Galois Theory}, version 5.10, September 2022.
\href{https://www.jmilne.org/math/CourseNotes/ft.html}{Author's notes}.

\bibitem{rz}
Luis Ribes and Pavel Zalesskii.
\emph{Profinite Groups}, second edition.
Ergebnisse der Mathematik und ihrer Grenzgebiete, volume 40.
Springer, 2010.

\bibitem{ln}
Rudolf Lidl and Harald Niederreiter.
\emph{Introduction to Finite Fields and Their Applications},
revised edition. Cambridge University Press, 1994.

\bibitem{stacks}
The Stacks Project Authors.
\emph{The Stacks Project}, Chapter 9, Fields, especially
\href{https://stacks.math.columbia.edu/tag/0BMI}{Section 9.22}.
Tags in the text identify the individual statements.

\bibitem{d07}
MathGenome.
\emph{Abstract Algebra} (D07), accepted notes and formalization,
particularly Chapters 12--14. MathGenome v0.21.0.

\bibitem{d10}
MathGenome.
\emph{Point-Set Topology: Unified Notes} (D10), Chapter~5.
Accepted revision included in MathGenome v0.22.0.
\end{thebibliography}
\end{document}
